\documentclass{report}
\usepackage{amsmath,amssymb}
\setlength{\parindent}{0mm}
\setlength{\parskip}{1em}
\begin{document}
\begin{center}
\rule{6in}{1pt} \
{\large Lior Horesh \\
{\bf Dynamics based geo-statistical prior sampling}}

IBM TJ Watson Research Center \\ 1101 Kitchawan Rd \\ Yorktown Heights \\ NY \\ 10598
\\
{\tt lhoresh@us.ibm.com}\\
Andrew Conn\\
Gijs van Essen\\
Eduardo Jimenez\\
Sippe Douma\\
Ulisses Mello\end{center}

Among the copious challenges inversion of large scale ill-posed problems
introduce lies quantification of uncertainty. Here the inherent
uncertainty of the problem is tackled in terms of prior sampling and our
focus is in attempting to realize how distinct are the model samples
considering their dynamic fingerprint. More specifically, we consider the
history matching inversion problem in the context of flow in porous
medium. Production data is collected from a sparse set of wells, and the
subsurface parameters are inferred through computationally intensive
inversion. Despite conscientious efforts to minimize the variability of
the solution space, the distribution of the posterior often remains
rather intrusive. In particular, geo-statisticians often propose large
sets of model prior samples that regardless of their apparent geological
distinction are almost entirely flow equivalent. As an antidote, in this
study a reduced space hierarchical clustering of flow-relevant indicators
is proposed for aggregation of these samples. Harnessing the dynamics as
dictated by the governing physics and the controls allow identifying
model characteristics that affect the dynamics. The effectiveness of the
method is demonstrated both with synthetic and real field data.


\end{document}
